
\documentclass{article}
\usepackage{colortbl}
\usepackage{makecell}
\usepackage{multirow}
\usepackage{supertabular}

\begin{document}

\newcounter{utterance}

\centering \large Interaction Transcript for game `cladder', experiment `full\_v1.5\_default', episode 5310 with refineonly{-}Qwen3.5{-}9B.
\vspace{24pt}

{ \footnotesize  \setcounter{utterance}{1}
\setlength{\tabcolsep}{0pt}
\begin{supertabular}{c@{$\;$}|p{.15\linewidth}@{}p{.15\linewidth}p{.15\linewidth}p{.15\linewidth}p{.15\linewidth}p{.15\linewidth}}
    \# & \multicolumn{2}{c}{Player} && \multicolumn{2}{c}{Game Master} \\
    \hline

    \theutterance \stepcounter{utterance}  
    & & & \multicolumn{4}{p{0.6\linewidth}}{
        \cellcolor[rgb]{0.9,0.9,0.9}{
            \makecell[{{p{\linewidth}}}]{
                \texttt{\tiny{[P1$\langle$GM]}}
                \texttt{You are an expert in causal inference. The following question is not a typical commonsense query, but rather a meticulously designed question created by a professor specializing in causal inference, intended to assess the students' mastery of the course content.} \\
\texttt{Imagine a self{-}contained, hypothetical world with only the following conditions, and without any unmentioned factors or causal relationships: Liking spicy food has a direct effect on occupation and salary. Skill has a direct effect on occupation and salary. Occupation has a direct effect on salary. Skill is unobserved. Method 1: We look at how liking spicy food correlates with salary case by case according to occupation. Method 2: We look directly at how liking spicy food correlates with salary in general. To understand how liking spicy food affects salary, is it more correct to use the Method 1 than Method 2?} \\
\texttt{Start your answer with 'Yes' or 'No', followed by additional reasoning or evidence to support your explanation.} \\
            }
        }
    }
     \\ \\

    \theutterance \stepcounter{utterance}  
    & \multicolumn{4}{p{0.6\linewidth}}{
        \cellcolor[rgb]{0.561,0.737,0.561}{
            \makecell[{{p{\linewidth}}}]{
                \texttt{\tiny{[P1$\rangle$GM]}}
                \texttt{No} \\
\\ 
\texttt{**Reasoning:**} \\
\\ 
\texttt{In this hypothetical world, the causal structure is defined by the following directed edges:} \\
\texttt{1.  **Liking Spicy Food** \$\textbackslash{}rightarrow\$ **Occupation**} \\
\texttt{2.  **Liking Spicy Food** \$\textbackslash{}rightarrow\$ **Salary** (Direct effect)} \\
\texttt{3.  **Skill** \$\textbackslash{}rightarrow\$ **Occupation**} \\
\texttt{4.  **Skill** \$\textbackslash{}rightarrow\$ **Salary**} \\
\texttt{5.  **Occupation** \$\textbackslash{}rightarrow\$ **Salary**} \\
\\ 
\texttt{Here, **Skill** is an unobserved confounder that affects both **Occupation** and **Salary**. **Occupation** acts as a mediator for the effect of **Liking Spicy Food** on **Salary**, but it is also a collider on the path between **Liking Spicy Food** and **Skill** (since both influence Occupation).} \\
\\ 
\texttt{Let's analyze the two methods:} \\
\\ 
\texttt{*   **Method 2 (Total Correlation):** This method examines the raw correlation between **Liking Spicy Food** and **Salary**. The observed association includes:} \\
\texttt{    1.  The direct causal effect (**Liking Spicy Food** \$\textbackslash{}rightarrow\$ **Salary**).} \\
\texttt{    2.  The indirect causal effect via the mediator (**Liking Spicy Food** \$\textbackslash{}rightarrow\$ **Occupation** \$\textbackslash{}rightarrow\$ **Salary**).} \\
\texttt{    3.  Spurious association due to the confounder **Skill**. Since **Skill** influences both **Occupation** and **Salary**, and **Liking Spicy Food** influences **Occupation**, there is a non{-}causal path: **Liking Spicy Food** \$\textbackslash{}rightarrow\$ **Occupation** \$\textbackslash{}leftarrow\$ **Skill** \$\textbackslash{}rightarrow\$ **Salary**.} \\
\texttt{    Because **Skill** is unobserved, Method 2 cannot distinguish the true causal effects from the confounding bias induced by **Skill**. It yields a biased estimate of the total causal effect.} \\
\\ 
\texttt{*   **Method 1 (Stratification by Occupation):** This method conditions on **Occupation**.} \\
\texttt{    *   Conditioning on **Occupation** blocks the indirect path through the mediator (**Liking Spicy Food** \$\textbackslash{}rightarrow\$ **Occupation** \$\textbackslash{}rightarrow\$ **Salary**). This allows us to isolate the **direct effect** of **Liking Spicy Food** on **Salary**.} \\
\texttt{    *   Crucially, conditioning on **Occupation** also blocks the backdoor path involving the confounder **Skill** (**Liking Spicy Food** \$\textbackslash{}rightarrow\$ **Occupation** \$\textbackslash{}leftarrow\$ **Skill** \$\textbackslash{}rightarrow\$ **Salary**). By fixing the occupation, we sever the link through which **Skill** creates a spurious correlation between **Liking Spicy Food** and **Salary**.} \\
\texttt{    *   Since **Skill** is the only unobserved confounder and it affects the outcome only through **Occupation** (and directly), and since we are} \\
            }
        }
    }
    & & \\ \\

    \theutterance \stepcounter{utterance}  
    & & & \multicolumn{2}{p{0.3\linewidth}}{
        \cellcolor[rgb]{0.9,0.9,0.9}{
            \makecell[{{p{\linewidth}}}]{
                \texttt{\tiny{[GM$|$GM]}}
                \texttt{No} \\
            }
        }
    }
    & & \\ \\

    \theutterance \stepcounter{utterance}  
    & & & \multicolumn{2}{p{0.3\linewidth}}{
        \cellcolor[rgb]{0.9,0.9,0.9}{
            \makecell[{{p{\linewidth}}}]{
                \texttt{\tiny{[GM$|$GM]}}
                \texttt{no} \\
            }
        }
    }
    & & \\ \\

    \theutterance \stepcounter{utterance}  
    & & & \multicolumn{2}{p{0.3\linewidth}}{
        \cellcolor[rgb]{0.9,0.9,0.9}{
            \makecell[{{p{\linewidth}}}]{
                \texttt{\tiny{[GM$|$GM]}}
                \texttt{game\_result = WIN} \\
            }
        }
    }
    & & \\ \\

\end{supertabular}
}

\end{document}
